% ===================================================================
%  BẢNG Số 13 — Khách sạn. --- Nhà hàng. Quán cà phê.
%
%  Traduction, faite en 2026, en vietnamien standard d'aujourd'hui,
%  sur l'ido de texto/io. Regles de la colonne : voir 10-bang-01.tex.
%
%  L'HOTEL EST UNE INSTITUTION IMPORTEE, ET SON PERSONNEL PORTE DEUX
%  SORTES DE NOMS. Ceux qui existaient deja se disent en vietnamien —
%  « người hầu phòng » la femme de chambre, « người nấu bếp » le
%  cuisinier, « người coi cửa » le concierge, « người đánh xe » le
%  cocher. Ceux que l'hotel a inventes gardent leur nom francais parce
%  qu'il est venu avec l'uniforme : « bồi », de \textit{boy}, pour le
%  garcon ; « gác-dông », de \textit{garcon}, encore vivant en 1926 ;
%  et « pooc-tê », de \textit{porteur}. On garde « bồi » seul, qui
%  est reste, et on rend les deux autres par le composé.
%
%  LES TROIS REPAS ONT ETE LE POINT DELICAT DU QUEBEC ; ILS NE LE
%  SONT PAS ICI. Le vietnamien compte les repas par le moment du jour
%  — « bữa sáng » le repas du matin, « bữa trưa » celui de midi,
%  « bữa tối » celui du soir — et aucun des trois ne peut glisser.
%  L'ido distingue lui-meme « dejuneto », le petit repas monte a la
%  chambre, et « dejuno », celui de la table d'hote : bữa sáng et bữa
%  trưa, sans qu'il y ait rien a trancher.
%
%  ET « TABLE D'HOTE » N'A PAS D'EQUIVALENT : on la rend par « bàn ăn
%  chung », la table ou l'on mange ensemble, qui dit exactement ce
%  qu'elle etait — un repas servi a heure dite, sans choix, pour tous
%  les pensionnaires a la fois.
%
%  LES JEUX DU CAFE SE PARTAGENT COMME CEUX DU TABLEAU 2. « Cờ » est
%  le mot general du jeu de plateau, et il sert de tete a trois des
%  cinq : cờ vua les echecs — le jeu du roi —, cờ đam les dames, cờ
%  tào cáo le trictrac. Les deux autres sont arrives par le port et
%  gardent leur nom : bi-a le billard, đô-mi-nô les dominos.
% ===================================================================

%%K t13-noto-1 noto
\VUnotes{143.2mm}{%
\textit{\VUgras{(*) Về chi tiết của căn phòng, xem bảng số 6.}}%
}

%%K t13-apar-1 apar
\VUsaut{3.0mm}
\VUtitre{15.0pt}{60}{BẢNG Số 13}
\VUsaut{1.6mm}
\VUpk{10.2pt}{40}{Khách sạn. --- Nhà hàng.}
\VUsaut{2.0mm}
\VUpk{10.2pt}{40}{Quán cà phê.}
\VUsaut{3.4mm}

%%K t13-01-1 p
1. --- Đến Roay-ăng chiều hôm qua, em trọ ở « \textit{Khách sạn Đại Dương} », nơi một
người bạn đã giới thiệu cho em.

%%K t13-01-2 p
Người ta cho em một \VUgras{căn phòng} \textsuperscript{(2)} đẹp ở \VUgras{tầng hai}
\textsuperscript{(1)}, nơi em đã ngủ rất ngon (*).

%%K t13-02-1 p
2. --- Sáng nay, khi ra khỏi phòng, nơi \VUgras{chị hầu phòng}
\textsuperscript{(3)} đã mang bữa sáng lên, em sang \VUgras{phòng hút thuốc}
\textsuperscript{(5)}, cũng dùng làm \VUgras{phòng đọc} \textsuperscript{(4)}, để xem
\VUgras{báo} \textsuperscript{(6)} và hút một điếu \VUgras{thuốc lá}
\textsuperscript{(8)}. Đọc xong, em xuống bằng \VUgras{thang máy}
\textsuperscript{(9)} và đi dạo trong thành phố.

%%K t13-03-1 p
3. --- Vì quên hỏi \VUgras{người tiếp tân} \textsuperscript{(12)} của khách sạn mấy giờ
thì dọn bữa ở \VUgras{bàn ăn chung} \textsuperscript{(11)}, em về sớm và nhân dịp ấy đi
tắm ở \VUgras{phòng tắm} \textsuperscript{(10)} của khách sạn. Rồi em trở lại
\VUgras{quầy} \textsuperscript{(15)} để gọi điện cho một người bạn. Nhưng \VUgras{máy
điện thoại} \textsuperscript{(13)} đang bận; thấy em lúng túng, \VUgras{chị thu ngân}
\textsuperscript{(14)}, người đang ghi một \VUgras{hóa đơn} \textsuperscript{(17)} vào
\VUgras{sổ khách trọ} \textsuperscript{(16)}, nhờ \VUgras{người coi cửa}
\textsuperscript{(18)} sai \VUgras{chú bồi} \textsuperscript{(19)} đi làm việc ấy giúp em
bằng \VUgras{xe đạp} \textsuperscript{(20)}.

%%K t13-04-1 p
4. --- Em cảm ơn chị và ra ngoài để cho tiền \VUgras{người khuân vác}
\textsuperscript{(24)} (người làm công việc nặng), người đã chở từ nhà ga về, trên chiếc
\VUgras{xe kéo tay} \textsuperscript{(25)} của mình, \VUgras{hộp mũ}
\textsuperscript{(26)}, \VUgras{va-li} \textsuperscript{(27)} và \VUgras{túi hành lý}
\textsuperscript{(28)} của em. \VUgras{Người đưa thư} \textsuperscript{(21)} vừa đến,
trao cho em một \VUgras{lá thư} \textsuperscript{(22)}.

%%K t13-05-1 p
5. --- Em còn nán lại một lúc trên bậc cửa, gần \VUgras{thùng thư}
\textsuperscript{(23)}, để nhìn cảnh đi lại ngoài phố. \VUgras{Người đánh xe}
\textsuperscript{(30)} của chiếc \VUgras{xe ngựa chở khách} \textsuperscript{(29)} đang
chất lên xe chiếc \VUgras{hòm} \textsuperscript{(31)} của một \VUgras{người khách}
\textsuperscript{(32)} mà \VUgras{ông chủ khách sạn} \textsuperscript{(33)} đang tiễn.
Một \VUgras{chú đưa điện tín} \textsuperscript{(35)} trẻ thong thả mang đến một
\VUgras{bức điện} \textsuperscript{(34)} cho một khách của khách sạn; một \VUgras{khách
du lịch} \textsuperscript{(36)}, đeo chéo vai chiếc \VUgras{máy ảnh}
\textsuperscript{(38)}, đang xem cuốn \VUgras{sách hướng dẫn} \textsuperscript{(37)} của
mình.

%%K t13-tit-1 sub
\VUsaut{4.0mm}
\VUpk{10.2pt}{40}{Giờ ăn trưa.}
\VUsaut{3.0mm}

%%K t13-06-1 p
6. --- Trước hàng rào sắt, một \VUgras{người chạy việc} \textsuperscript{(39)} thản
nhiên ngủ gà ngủ gật giữa mấy cái \VUgras{móc gánh} \textsuperscript{(40)} và cái
\VUgras{hộp đánh giày} \textsuperscript{(41)} của mình, trong khi một \VUgras{cô thợ
giặt} \textsuperscript{(42)} trẻ, mang một \VUgras{giỏ} \textsuperscript{(43)} đồ vải
lớn, cười vẻ mặt trịnh trọng của \VUgras{người quản lý nhà ăn}
\textsuperscript{(44)} đang đứng gần cửa.

%%K t13-07-1 p
7. --- Trông thấy \VUgras{người nấu bếp} \textsuperscript{(45)}, người vừa ra khỏi
\VUgras{bếp} \textsuperscript{(47)} của mình ở \VUgras{tầng hầm} \textsuperscript{(48)}
để sai một \VUgras{chú phụ bếp} \textsuperscript{(46)} đi làm việc vặt, em nhớ ra rằng
giờ ăn trưa sắp đến, và em vào nhà ăn, băng qua cái sân nơi một \VUgras{người lái xe}
\textsuperscript{(50)} đang cho \VUgras{ô-tô} \textsuperscript{(49)} của mình vào chỗ đỗ,
trong khi một \VUgras{thợ máy} \textsuperscript{(52)} sửa, theo lời chỉ dẫn của
\VUgras{người đi xe} \textsuperscript{(53)}, một chiếc \VUgras{xe ba bánh chạy dầu}
\textsuperscript{(51)} vừa gặp tai nạn.

%%K t13-08-1 p
8. --- Em hỏi một \VUgras{chú bồi} \textsuperscript{(54)} nhỏ đang bưng những \VUgras{cốc
bia} \textsuperscript{(55)} xem bàn ăn chung có ở dưới hiên không, và vì chú trả lời có,
em chọn chỗ ở một trong các bàn và gọi một bữa ăn tuyệt hảo.

%%K t13-noto-2 noto
\VUnotes{143.2mm}{%
\textit{\VUgras{(*) Nhờ danh sách các món ăn có trong từ vựng, thầy giáo sẽ dễ dàng thay
đổi thực đơn dưới đây.}}%
}

%%K t13-tit-2 sub
\VUsaut{4.0mm}
\VUpk{10.2pt}{40}{Một bữa trưa tuyệt hảo.}
\VUsaut{3.0mm}

%%K t13-09-1 p
9. --- Người bồi mang đến cho em \VUgras{thực đơn} (*) bữa trưa và danh sách rượu vang.
Em chọn và gọi làm \VUgras{món khai vị} món \VUgras{tôm} với \VUgras{bơ}, rồi làm món
đầu, \VUgras{lòng bò hầm kiểu Caen}. Em không ăn \VUgras{cá}, nhưng gọi một phần
\VUgras{đùi} cừu non và, làm \VUgras{món rau}, \VUgras{rau bina} nấu kem. Sau đó, vì
không món nào trong các \VUgras{món ngọt} hợp ý em, em bảo mang lên nhiều thứ tráng
miệng khác nhau : \VUgras{phô mai}, \VUgras{trái cây} và \VUgras{bánh quy}. Em còn thêm
một chai \VUgras{rượu vang} Saint-Émilion tuyệt hảo và, rất hài lòng về bữa ăn, em rời
bàn sau khi cho \VUgras{người bồi} \textsuperscript{(57)} một khoản \VUgras{tiền thưởng}
\textsuperscript{(59)} hậu hĩnh.

%%K t13-10-1 p
10. --- Thấy em sắp đi, \VUgras{người quản lý} \textsuperscript{(60)} mời em lên ban
công để ngắm quang cảnh tuyệt vời của \VUgras{Đại dương} \textsuperscript{(62)}. Xa xa
thấy những chiếc \VUgras{thuyền} \textsuperscript{(63)} đánh cá nhỏ, những chiếc
\VUgras{thuyền buồm} \textsuperscript{(64)} và một chiếc \VUgras{tàu khách}
\textsuperscript{(65)} khổng lồ mà bóng đen nổi bật trên nền \VUgras{mây}
\textsuperscript{(67)} trắng ở \VUgras{chân trời} \textsuperscript{(66)}. Một làn gió nhẹ
làm bay \VUgras{lá cờ} \textsuperscript{(70)} cắm trên \VUgras{biển hiệu}
\textsuperscript{(69)} của \VUgras{sảnh} \textsuperscript{(68)}.

%%K t13-tit-3 sub
\VUsaut{3.0mm}
\VUpk{10.2pt}{40}{Giải trí.}
\VUsaut{3.0mm}

%%K t13-11-1 p
11. --- Cảnh ấy thú vị đến nỗi em quyết định ngày mai sẽ lên sân thượng của quán cà phê,
mang theo một chiếc \VUgras{ống nhòm} \textsuperscript{(71)} hoặc một chiếc \VUgras{kính
viễn vọng} \textsuperscript{(72)}.

%%K t13-12-1 p
12. --- Rời ban công, em sang quán cà phê của khách sạn để uống một \VUgras{thức uống}
\textsuperscript{(73)} và chơi một \VUgras{ván bi-a} \textsuperscript{(75)}. Em đi thẳng
qua phòng cà phê, nơi nhiều khách đang chơi \VUgras{cờ vua} \textsuperscript{(78)},
\VUgras{cờ đam} \textsuperscript{(79)}, \VUgras{đô-mi-nô} \textsuperscript{(80)},
\VUgras{cờ tào cáo} \textsuperscript{(81)} hoặc \VUgras{bài} \textsuperscript{(82)}, rồi
lên thẳng \VUgras{phòng bi-a} \textsuperscript{(74)}.

%%K t13-13-1 p
13. --- Khi ván đấu xong, em xuống tầng trệt và xin người bồi một chiếc \VUgras{phong bì}
\textsuperscript{(84)}, ít \VUgras{giấy} \textsuperscript{(83)}, một con \VUgras{tem}
\textsuperscript{(85)}, một cái \VUgras{ngòi bút} \textsuperscript{(87)} và ít
\VUgras{mực} \textsuperscript{(86)} để viết thư.

%%K t13-13-2 p
Rồi em gọi một \VUgras{người đánh xe} \textsuperscript{(92)} có chiếc \VUgras{xe}
\textsuperscript{(88)} đỗ trước quán cà phê. Em hỏi giá theo giờ hay theo cuốc và, khi
thỏa thuận xong, ông mở \VUgras{cửa xe} \textsuperscript{(89)} cho em và mời em lên. Ông
trèo lại lên \VUgras{ghế ngồi} \textsuperscript{(91)} của mình, cho ngựa chạy bằng một
cái vung nhẹ chiếc \VUgras{roi} \textsuperscript{(93)}, và chúng em lên đường cho một
chuyến dạo chơi thú vị.
\VUsaut{6.0mm}
\VUfilet{22mm}
